% !TEX root = Tesi_Corsi_Leonardo_656276.tex
\chapter{Strumenti utilizzati}
\label{ch:strumenti}

Il progetto si articola in tre viste dello stesso sistema: il simulatore
in C, i programmi in assembly che vi girano sopra, e il modello del
circuito in SystemVerilog. Un'architettura di riferimento comune le
tiene insieme.

\section{RISC-V e il sottoinsieme RV32I}
\label{sec:riscv}

RISC-V è una architettura dell'insieme di istruzioni (\emph{Instruction
Set Architecture}, ISA) aperta, organizzata come un nucleo base
obbligatorio e un insieme di estensioni opzionali~\cite{RiscvSpec}. Il
nucleo intero a 32 bit, RV32I, conta poche decine di istruzioni: \emph{load} e \emph{store},
operazioni aritmetico-logiche fra registri o con un immediato, salti
condizionati e incondizionati, e le due istruzioni che costruiscono le
costanti a 32 bit. Non ha moltiplicazione né divisione: queste operazioni
sono nell'estensione M, che la convenzione di nomenclatura di RISC-V
aggiunge come suffisso al nome del nucleo: RV32I con l'estensione M si
indica RV32IM.

Per un array di processori, RISC-V risponde a tre esigenze:

\begin{itemize}
    \item \emph{Dimensione.} Una cella deve essere piccola, perché di
          celle ce ne sono $R \times C$, e RV32I è fra le ISA complete
          più piccole: solo RV32E, che dimezza i registri, lo è di più.
    \item \emph{Apertura.} L'architettura è documentata nel dettaglio
          e priva di vincoli d'uso.
    \item \emph{Estendibilità.} RISC-V riserva esplicitamente parte
          dello spazio di codifica alle estensioni non standard: si
          possono aggiungere istruzioni proprie \emph{dentro}
          l'architettura, invece che accanto a essa.
\end{itemize}
I programmi di prova sono compilati per \texttt{rv32i} puro: è il valore
di default della variabile \texttt{ARCH} nel \emph{Makefile}. Il vincolo
è intenzionale: se l'interprete non implementa un'istruzione,
l'assemblatore deve rifiutarla prima ancora che il programma venga
eseguito. L'unica eccezione è il programma di moltiplicazione di
matrici, che ha bisogno dell'istruzione \texttt{mul}: la regola di
compilazione del suo solo file oggetto sovrascrive \texttt{ARCH} a
\texttt{rv32im}, mentre il default \texttt{rv32i} resta invariato per
tutti gli altri bersagli di compilazione.

\subsection{Lo spazio di codifica \emph{custom-0}}
\label{sec:custom0}

Le quattro istruzioni di comunicazione usano l'opcode \texttt{custom-0}
(\texttt{0x0B}), uno di quelli che la specifica lascia liberi per le estensioni
non standard, con il formato I-type:

\begin{center}
\begin{tabular}{ccccc}
    \toprule
    \texttt{imm[11:0]} & \texttt{rs1} & \texttt{funct3} & \texttt{rd} & \texttt{opcode} \\
    31\dots20 & 19\dots15 & 14\dots12 & 11\dots7 & 6\dots0 \\
    \midrule
    direzione & sorgente & operazione & destinazione & \texttt{0x0B} \\
    \bottomrule
\end{tabular}
\end{center}
La direzione sta nei due bit bassi dell'immediato ($0 = $ nord, $1 = $ est,
$2 = $ sud, $3 = $ ovest) e il campo \texttt{funct3} distingue le quattro
operazioni. Il dettaglio della semantica è nella
sezione~\ref{sec:istruzioni}.

Per emetterle non è stato necessario modificare l'assemblatore. La direttiva
\texttt{.insn} di GNU Binutils~\cite{Binutils} esiste proprio per gli opcode
custom e permette di scrivere una istruzione dando i campi di codifica;
avvolta in una macro, si usa poi come una istruzione qualsiasi:

\begin{lstlisting}[caption={La macro con cui il programma della cella emette \texttt{ISRDY}},label={lst:macro},language=RISCV,float=htbp]
.macro ISRDY rd, dir
    .insn i 0x0B, 0x2, \rd, x0, \dir
.endm
\end{lstlisting}
La toolchain resta quella ufficiale e non modificata: non è stato necessario
ricompilare l'assemblatore per aggiungere supporto alle quattro istruzioni, e
i file assembly restano assemblabili da chiunque disponga della toolchain
RISC-V standard.

\section{C, assembly e il sistema di build}
\label{sec:c-assembly}

Il simulatore è scritto in C, come l'interprete a singolo processore da
cui il progetto è partito~\cite{TesiStella}. Il canale è un oggetto di
poche decine di byte di cui interessa il \emph{layout} esatto, e il
ciclo di simulazione lo visita milioni di volte: è realizzato
interamente con funzioni \texttt{static inline} in un solo file di
intestazione, così il compilatore le espande nel punto di chiamata e il
costo dell'astrazione sparisce.

I programmi che girano sull'array sono in assembly RISC-V: le istruzioni
di comunicazione non esistono in nessun linguaggio di livello superiore, e
la logica di sincronizzazione (attendere che un canale sia pronto,
ritentare una pubblicazione rifiutata) è quindi scritta nel programma
assembly stesso, non nell'interprete, che si limita a eseguire l'istruzione
e a restituirne l'esito.

Il tutto è gestito da un \emph{Makefile} che compila l'interprete,
assembla e collega i programmi in file ELF, li carica nella griglia, esegue la
suite di verifica e rigenera i dati delle misure. Le prove non sono esecuzioni
singole ma un \emph{parameter sweeping} (da qui in avanti semplicemente \emph{sweep}): 
ogni programma viene eseguito una volta per ciascuna combinazione dei parametri che
lo riguardano, e ogni esecuzione stampa una riga in formato CSV. 
La composizione dello sweep è nella sezione~\ref{sec:metodologia}. I file ELF sono prodotti dalla
catena \texttt{riscv64-unknown-elf} configurata per \texttt{rv32i}, e caricati
dal simulatore con un lettore ELF minimale.

\section{SystemVerilog e la catena di strumenti hardware}
\label{sec:systemverilog}

Il protocollo è nato come modello in C, ma il suo scopo è descrivere un
circuito. Il canale e l'interfaccia di comunicazione di una cella sono
stati perciò riscritti in SystemVerilog, per misurare il costo
dell'interfaccia in termini di flip-flop e di porte logiche.

La catena usata è interamente libera e si compone di tre strumenti con tre ruoli
distinti:

\begin{itemize}
    \item \emph{Icarus Verilog} esegue i due \emph{testbench}, che ripercorrono
          sul modello RTL gli stessi casi limite già verificati sul modello C;
    \item \emph{Verilator}, in modalità \texttt{-{}-lint-only}, controlla la
          conformità del codice. È più severo di Icarus, che accetta in silenzio
          costrutti che uno strumento di sintesi rifiuterebbe;
    \item \emph{Yosys} sintetizza il canale e ne riporta la statistica di
          celle. È lo strumento utilizzato per i numeri della sezione~\ref{sec:rtl}.
\end{itemize}

\section{OpenMP}
\label{sec:openmp}

Il simulatore fa avanzare $R \times C$ processori che per costruzione non
condividono memoria, e li fa avanzare tutti dello stesso passo. È il caso
tipico di parallelismo di dati, e OpenMP è lo strumento adatto proprio perché
richiede pochissimi interventi per parallelizzare il codice del simulatore: 
le due fasi del ciclo di clock sono due cicli \texttt{for} sull'array delle celle, 
e diventano paralleli con la semplice introduzione di una direttiva \texttt{\#pragma 
omp parallel for} per ciascuno dei due for.

La direttiva non richiede alcun meccanismo di sincronizzazione
aggiuntivo: nessun \emph{lock}, nessuna sezione critica, nessuna
riduzione. La ragione sta nella struttura del ciclo di clock ed è
discussa nella sezione~\ref{sec:parallelizzazione}: la barriera implicita
a fine ciclo parallelo, che OpenMP fornisce di serie, coincide
esattamente con la separazione fra le due fasi.

La misura del guadagno, e delle griglie su cui il parallelismo costa invece di
rendere, è nella sezione~\ref{sec:prestazioni}.